% 后置内容：参考文献、创新成果、声明和致谢。
% 由main.tex加载，不单独编译；排版设置仍在geshi.tex。
\printbibliography[heading=bibintoc,title={参考文献}]
% ==================== 07 自动参考文献 ====================
% ==================== 08 创新成果 ====================
\achievementtitle
\achievementcategory{（一）发表的学术论文}
\begin{achievementlist}
\item 作者. 论文题目[J]. 期刊名称, 年, 卷(期): 起止页码.（对应本论文第1章）
\end{achievementlist}
\achievementcategory{（二）申请及已获得的专利}
\begin{achievementlist}
\item 发明人. 专利名称: 国别, 专利号[P]. 公开日期.
\end{achievementlist}
\achievementcategory{（三）参与的科研项目及获奖情况}
\begin{achievementlist}
\item 项目参与人. 项目名称, 项目来源, 课题编号.
\end{achievementlist}
% 没有的成果类别连同其列表一起删除；成果不得混入正文参考文献。

% ==================== 09 原创性声明与使用权限 ====================
\begin{thesisstatement}
\statementheading{学位论文原创性声明}
本人郑重声明：此处所提交的学位论文《\thesisTitle》，是本人在导师指导下，在哈尔滨工业大学攻读学位期间独立进行学术研究工作或专业实践工作所取得的成果，且学位论文中除已标注引用文献、资料的部分外不包含他人完成或已发表的成果；对本学位论文的学术研究工作或专业实践工作做出重要贡献的个人和集体，均已在文中以明确方式注明；对使用AI工具的情况均已在文中以明确方式标注。

\originalitysignature
\permissionheading
学位论文是研究生在哈尔滨工业大学攻读学位期间完成的成果，知识产权归属哈尔滨工业大学。学位论文的使用权限如下：
\begin{permissionclauses}
\item 学校可以采用影印、缩印或其他复制手段保存研究生上交的学位论文，并向国家图书馆报送学位论文；
\item 学校可以将学位论文部分或全部内容编入有关数据库进行检索和提供相应阅览服务；
\item 研究生毕业后发表与此学位论文研究成果相关的学术论文和其他成果时，应征得导师同意，且第一署名单位为哈尔滨工业大学。
\end{permissionclauses}
保密论文在保密期内遵守有关保密规定，解密后适用于此使用权限规定。

本人知悉学位论文的使用权限，并将遵守有关规定。

\permissionsignatures
\end{thesisstatement}

% ==================== 10 致谢 ====================
\section*{致\quad 谢}
在此填写致谢。

